\section{Delooping with torsors}
\label{torsor-delooping}
In this section, we recall the other classical approach to constructing
deloopings of groups by using $G$-torsors, which originates in classical
constructions of algebraic topology~\cite{demazure1970groupes}. Most of the
material of the section is already known, for which reason proofs are not much
detailed. A more in-depth presentation can be found in recent works such
as~\cite{symmetry,buchholtz2023central}.

\paragraph{Group actions}
% \label{group-action}
Given a group~$G$ and a set~$A$, an \emph{action} of $G$ on $A$ is a group
morphism from $G$ to $A \simeq A$, that is a map $\alpha : G\to (A \simeq A)$
such that
\begin{align}
  \label{action-equations}
  \alpha(xy)&= \alpha(x) \circ \alpha(y)
  &
  \alpha(1) &= \id_A
\end{align}
% $\alpha(xy)= \alpha(x) \circ \alpha(y)$
% and
% $\alpha(1) = \id_A$
for all $x,y:G$.

A \emph{$G$-set} is a set equipped with an action of~$G$, and we write $\GSet$
for the type of $G$-sets. We often simply denote a $G$-set by the associated
action $\alpha$ and write $\dom(\alpha)$ for the set on which~$G$ acts.

\begin{lemma}[\agda{GSetProperties}{isGroupoidGSet}]
  \label{gset-groupoid}
  The type $\GSet$ is a groupoid.
\end{lemma}
\begin{proof}
  The type of sets is a groupoid~\cite[Theorem 7.1.11]{hottbook}. Given a set
  $A$, the type of functions $A\to A$ is a set~\cite[Theorem 7.1.9]{hottbook}
  and thus a groupoid. Finally,
  % the equations \cref{action-equations}
  the axioms~\cref{action-equations} of actions are propositions (because $A$
  is a set) and thus groupoids. We conclude since groupoids are closed under
  $\Sigma$-types~\cite[Theorem 7.1.8]{hottbook}.
\end{proof}

Given $G$-sets $\alpha$ and $\beta$, a \emph{morphism} between them consists of
a function $f:\dom{\alpha}\to \dom{\beta}$ which preserves the group action, in
the sense that for every $x:G$ and $a:\dom{\alpha}$, we have
\begin{equation}
  \label{G-set-morphism}
  \beta(x)(f(a))=f(\alpha(x)(a))
  \text.
\end{equation}
A morphism which is also an equivalence is called an \emph{isomorphism} and we
write $\GSetEquiv\alpha\beta$ for the type of isomorphisms between $\alpha$ and
$\beta$. We write $\Aut(\alpha)$ for the type of automorphisms
$\GSetEquiv\alpha\alpha$, which is a group under composition.
%
% \subsubsection{Identities of~$G$-sets}
The equalities between $G$-sets can be conveniently characterized as follows.

\begin{proposition}[\agda{GSetProperties}{GSet≡Decomp}]
  \label{Gset-equalities}
  Given two $G$-sets $\alpha$ and $\beta$, an equality between them consists of
  an equality $p:\dom{\alpha}=\dom{\beta}$ such that the function induced by
  transport along~$p$, namely $\transport p:\dom\alpha\to\dom\beta$, is a
  morphism of $G$-sets.
\end{proposition}
\begin{proof}
  The characterization of equalities between
  $\Sigma$-types~\cite[Theorem~2.7.2]{hottbook} entails that an equality between
  $(\dom{\alpha},\alpha)$ and $(\dom{\beta},\beta)$ is a pair consisting of an
  equality $p:\dom{\alpha}=\dom{\beta}$ and an equality $q:\transport{p}(\alpha)=\beta$ (we can forget
  about the equality between the components expressing the properties required
  for group actions since those are propositions). By
  \cite[Lemma~2.9.6]{hottbook} and function extensionality, we finally have that
  the type of~$q$ is equivalent to the type
  $\beta(x)\circ\transport p=\transport p\circ\alpha(x)$.
\end{proof}

\noindent
It easily follows from this proposition that any equality between $G$-sets
induces an isomorphism of $G$-sets, as customary for equalities between
algebraic structures~\cite[Section~2.14]{hottbook}. In fact, this map from
equalities to isomorphisms can itself be shown to be an equivalence:

\begin{proposition}[\agda{GSetProperties}{GSetPath}]
  \label{GSet-SIP}
  Given $G$-sets $\alpha$ and $\beta$, the canonical function
  \[
    (\alpha=\beta)\to(\GSetEquiv\alpha\beta)
  \]
  is an equivalence.
  %
  Moreover, given a $G$-set $\alpha$, the induced equivalence
  \[
    (\alpha=\alpha)\equivto(\GSetEquiv\alpha\alpha)
  \]
  is compatible with the canonical group structures on both types.
\end{proposition}
\begin{proof}
  % Finally, we can show that $G$-Sets respect the Structure Identity Principle,
  % \ie that isomorphic $G$-Sets can be identified. A proof is given in the agda
  % code using the fundamental theorem of identity types. However, the cubical
  % library offers a streamlined way of applying this principle to various
  % algebraic structures, following Escardo's work on the
  % topic~\cite{escardo2019introduction}, which may and should apply here.
  %
  This is actually an instance of a more general correspondence between
  equalities and isomorphisms of algebraic structures, which is known under the
  name of \emph{structure identity
    principle}, see~\cite{coquand2013isomorphism} and \cite[Section~9.8]{hottbook}, and can be understood as a
  generalisation of univalence for types having an algebraic structure.
  % \AR{This is a rather one-sided description of the history of the structure
  % identity principle.  For example the HoTT Book already contains (in §9.8) a
  % proof of an SIP for some class of structures. The paper Coquand, Danielsson
  % Isomorphism is equality \url{http://dx.doi.org/10.1016/j.indag.2013.09.002}
  % \cite{coquand2013isomorphism} formalizes an SIP in Agda and predates [5].  I
  % won't list out every paper about the SIP here, but it's worth noting that
  % different people have formulated SIP's in different ways and in different
  % generality.}
\end{proof}

\paragraph{Torsors}
For any group~$G$, there is a canonical $G$-set called the \emph{principal $G$-torsor} and
noted $\P G$, corresponding to the action of $G$ on itself by left
multiplication. Moreover, its group of automorphisms is precisely the group~$G$:

\begin{proposition}[\agda{Deloopings}{PGloops}]
  \label{loop-PG}
  Given a group~$G$, we have an equality of groups
  \[
    (\GSetEquiv{\P G}{\P G})=G
  \]
\end{proposition}
\begin{proof}
  % Note that any element $f:\Aut\P G$, \ie morphism of actions $f:\P G\to\P G$,
  % satisfies $xf(y)=f(xy)$ for any $x,y:G$.
  The two functions
  % \begin{align*}
    % \phi:\Aut\P G&\to G
    % &
    % \psi:G&\to\Aut\P G
    % \\
    % f&\mapsto f(1)
    % &
    % x&\mapsto y\mapsto yx
    % \end{align*}
  $\phi:\Aut\P G\to G$ and $\psi:G\to\Aut\P G$, respectively defined by
  $\phi(f)\defd f(1)$ and $\psi(x)(y)\defd yx$ are mutually inverse group
  morphisms, see \cref{omitted-proofs}.
\end{proof}

The type $\GSet$ is thus ``almost'' a delooping of~$G$. Namely, it is a groupoid
(\cref{gset-groupoid}), which is pointed by $\P G$ and satisfies~$\Loop\GSet = G$ by
\cref{GSet-SIP,loop-PG}. It only lacks being connected, which can easily be
addressed. Given a pointed type~$A$, its \emph{connected component} $\Comp A$ is
the type $\Sigma(x:A).\ptrunc{\pt{}=x}$. It is well-known that this type is
pointed by $(\pt{},\ptrunq{\refl{}})$, connected and has the same loop space as
the original type, \ie we have $\Loop\Comp A=\Loop A$, see
\cref{connected-component}. We thus have:

\begin{theorem}[\agda{Deloopings}{torsorDeloops}]
  \label{PG-delooping}
  The connected component of~$\P G$ in $\GSet$, \ie the type
  $\Comp(\GSet,\P G)$, is a delooping of $G$.
\end{theorem}

\noindent
The elements of the connected component of the principal $G$-set are usually
called \emph{$G$-torsors}.
% Local Variables:
% TeX-master: "article"
% End:
